\paragraph{\S105, continuation. Abelian pairings, algebraic topoi and an unfinished universal property (pp.\ 399--406).}
This continuation begins with Case~2, introduced on p.\ 398. For $B=A^o$, (49) identifies $B^\wedge_{ab}$ with the category of additive cosheaves on $A^\wedge_{ab}$ valued in $(Ab)$. The corresponding pairing (50) commutes with small direct limits in each argument, and its symmetry supplies the analogous description of $A^\wedge_{ab}$. The equivalence (51) reduces such pairings to their values on $A\times B$. Formula (52) gives the free abelian group on $\mathrm{Hom}(b,a)$ when the arguments are the abelianizations of representables; (53) extends this to $\underline Z^{(F*G)}$ for presheaves $F,G$. The insertion on p.\ 400 explicitly says that \emph{both members}, rather than only the left side, commute with the relevant direct limits (pp.\ 399--400).

Remark~1 abstracts this situation to two additive categories $\mathcal P,\mathcal Q$ with a pairing to $(Ab)$, cocontinuous in each variable, which identifies each category with the additive cosheaves on the other. The presheaf examples are abelian topoi: abelian categories with exact small filtering direct limits and a small generating subcategory. The visible revisions add the filtering qualification and clarify that this terminology does not mean a category which is simultaneously an ordinary topos and an abelian category. Replacing abelian groups by modules over a commutative ring gives the corresponding $k$-linear examples. For a noncommutative ring, the two sides use right and left modules; the proposed pairing is the abelian-group pairing followed by quotienting the resulting bimodule by the differences $sx-xs$. Grothendieck explicitly says that he has not checked that this construction gives the desired duality. He does verify the ordinary one-object example through tensor product, evaluation at the regular module, and the associated tensor--Hom equivalences. Their comparison is an additive analogue of (39) when the target category has both kinds of small limits (pp.\ 400--402).

Remark~2 singles out the existence of a small generating family of ultraprojective objects. For a module category, one ultraprojective generator suffices; ultraprojectivity is decomposed into projectivity and finite presentation. The proposed general structure theorem describes the resulting “algebraic” abelian topoi as categories of additive functors from $P^o$ to $(Ab)$, for a small additive category $P$. One can take $P$ essentially small and canonically equal to the full subcategory of ultraprojective objects. For $A^\wedge_{ab}$, this is the abelian Karoubi envelope of $A$, or alternatively one may use $\mathrm{Add}(A)$ for more explicit computation; the text refers back to \S\S93 and 99 (p.\ 402). The notation $P^{\&}$ and the inclusion $\varepsilon_P:P\to P^{\&}$ lead to the expected equivalence (**) between cocontinuous additive functors out of $P^{\&}$ and additive functors out of $P$. The suggested proof imitates the nonadditive representable-density argument. Its dual gives the inverse-limit version and then (***), comparing sheaves and cosheaves on the dual algebraic abelian topoi. These statements are presented as the expected additive theory, with the proof still to be checked (p.\ 403).

The next page calls an abelian topos \emph{reflexive} when it belongs to such a dual pair. Algebraicity appears sufficient, and the author wonders whether it is necessary, at least for the stronger equivalence with arbitrary additive target categories and suitable compatibility as the target varies. Remark~3 records the central limitation: for a general dual pair he does not see how to construct the comparison functor from the pairing alone. The same perplexity occurs in the nonabelian case, so the algebraic examples should not be read as an established general duality theorem (p.\ 404). Remark~4 returns to the auxiliary category $\mathcal E(A)$ from p.\ 398. Inside the functors from $B^\wedge$ to sets, it is represented by the union of the cocontinuous functors and the representable functors. Their intersection is the Karoubi envelope of $A$, equivalently the opposite of the ultraprojective objects of $B^\wedge$. The margin note describes the union by commutation with either type of small limits. These constructions are then said to extend to a small additive $P$ (pp.\ 404--405).

Case~3 resumes the intended application of (39). With an additive target $\mathcal M$, initially under the extra pseudotopos assumptions on $\mathcal M$ and its opposite, (12) and (13) yield (54), relating $\mathcal M$-valued cosheaves on $A^\wedge_{ab}$ to sheaves on $B^\wedge_{ab}$. The earlier discussion suggests that those extra assumptions should be unnecessary. The required assertion (55) says that precomposition with abelianization gives the expected equivalence for any additive $\mathcal M$ with small direct limits. Formula (56) reduces it to additive functors on $\mathrm{Add}(A)$, and (57) is the inverse-limit dual. The four-arrow diagram (58) shows that the two left comparisons are equivalent conditions because the two right comparisons are already equivalences. Finally, identifying $A^\wedge_{ab}$ with $P^{\&}$ for $P=\mathrm{Add}(A)$ reduces the remaining task to (**). The author prefers to prove that general statement and ends by saying that the suggested proof should be checked during later scratchwork. The scan's handwritten reference in the insertion is (56), and its final wording is “proof given” followed by the correction from “there” to “here” (pp.\ 405--406).

